\chapter{Variáveis, Tipos e Operadores}

Todo programa existe para representar informações e operar sobre elas. Se você já teve contato com Python, provavelmente está acostumado a pensar em uma variável simplesmente como "uma caixa que guarda um valor" — você escreve idade = 20 e pronto, o Python descobre sozinho que aquilo é um número inteiro.

Em C, esse raciocínio precisa ser mais preciso. Antes de guardar qualquer valor, precisamos responder três perguntas:

\begin{itemize}
    \item Qual é o tipo desse dado?
    \item Quanto espaço ele ocupa na memória?
    \item Como esse valor é efetivamente representado, em bits, dentro desse espaço?
\end{itemize}

Essas perguntas podem parecer burocráticas no início, mas são exatamente o que torna C uma linguagem tão próxima do funcionamento real do computador. Neste capítulo, vamos construir a base para pensar em variáveis não como "caixas mágicas", mas como regiões de memória com um tipo bem definido.

\section{Variáveis e Declaração}

Uma variável é um espaço nomeado na memória do computador, usado para armazenar um valor que pode ser lido e modificado durante a execução do programa. Em C, toda variável precisa ser \textbf{declarada} antes de ser usada — ou seja, você precisa informar ao compilador qual é o tipo dela e qual nome ela vai ter.

\begin{lstlisting}
int idade;
idade = 20;

int ano = 2026;
\end{lstlisting}

Note que existem duas ações diferentes acontecendo:

\begin{itemize}

    \item \textbf{Declaração:} \texttt{int idade;} — aqui você está dizendo ao compilador "reserve espaço na memória para um número inteiro chamado idade".
    \item \textbf{Atribuição:} \texttt{idade = 20;} — aqui você está colocando o valor 20 nesse espaço já reservado.

\end{itemize}

Quando as duas coisas acontecem na mesma linha, como em \texttt{int ano = 2026;}, chamamos isso de inicialização — declaração e atribuição em um único passo. \\

\textbf{Regras para identificadores (nomes de variáveis) em C:}

\begin{enumerate}
    \item Podem conter letras, dígitos e underscore (\_).
    \item Não podem começar com um dígito.
    \item São case-sensitive (idade e Idade são variáveis diferentes).
    \item Não podem ser palavras reservadas da linguagem (como int, return, if).
\end{enumerate}

\textbf{A diferença em relação ao Python:} 
em Python, uma mesma variável pode receber valores de tipos diferentes ao longo do programa (\texttt{x = 10} e depois \texttt{x = "texto"} são ambos válidos). Em C isso não é possível — uma vez declarada com um tipo, a variável só pode armazenar valores compatíveis com aquele tipo durante toda a sua existência.


\begin{errocomum}

Declarar uma variável local e tentar usá-la antes de atribuir um valor:

\begin{lstlisting}
int total;
printf("%d", total); // comportamento indefinido!

\end{lstlisting}

Diferente de variáveis globais (que começam zeradas automaticamente), variáveis locais não têm um valor inicial garantido. Elas contêm "lixo de memória" — o que quer que estivesse naquele espaço antes. 

\textbf{Sempre} inicialize suas variáveis antes de usá-las.

    
\end{errocomum}



\section{Tipos de Dados Fundamentais}
C possui alguns tipos básicos, chamados de \textbf{tipos primitivos}, que servem de base para tudo o que vamos construir daqui em diante:

\begin{center}
\begin{tabular}{|l|l|l|}
\hline
\textbf{Tipo} & \textbf{Representa} & \textbf{Exemplo} \\
\hline
\texttt{char} & um único caractere & \texttt{char letra = 'A';} \\
\hline
\texttt{int} & um número inteiro & \texttt{int quantidade = 42;} \\
\hline
\texttt{float} & número real (precisão simples) & \texttt{float altura = 1.75f;} \\
\hline
\texttt{double} & número real (precisão dupla) & \texttt{double pi = 3.14159265;} \\
\hline
\texttt{void} & ausência de tipo/valor & usado principalmente em funções \\
\hline
\end{tabular}
\end{center}

Mais adiante, vamos discutir sobre o tipo \texttt{bool}, que representa valores verdadeiro/falso, em C com mais calma quando chegarmos às estruturas condicionais.

Por enquanto, não vamos nos aprofundar em modificadores como \texttt{short}, \texttt{long}, \texttt{signed} e \texttt{unsigned} — eles merecem atenção própria e voltaremos a eles na Seção~\ref{sec:constantes-modificadores}.

A ideia central desta seção é a seguinte:

\begin{quote}
\textit{Em C, o tipo faz parte da definição da variável e determina como seus bits devem ser interpretados.}
\end{quote}

Ou seja: os mesmos 4 bytes na memória podem representar um número inteiro ou um número em ponto flutuante, dependendo exclusivamente do tipo declarado. O tipo não é um detalhe burocrático — é o que dá significado aos bits armazenados.

\section{Representação dos Dados na Memória}
Este é um dos pontos em que C se diferencia claramente de linguagens de mais alto nível: podemos — e devemos — pensar em como os dados são fisicamente armazenados.

Considere a declaração:

\begin{lstlisting}
int x = 10;
\end{lstlisting}

Na maioria dos sistemas atuais, um \texttt{int} ocupa \textbf{4 bytes} (32 bits) de memória. O valor \texttt{10}, representado em binário, fica armazenado assim:

\begin{center}
\texttt{
\begin{tabular}{|c|c|c|c|}
\hline
00000000 & 00000000 & 00000000 & 00001010 \\
\hline
\end{tabular}
}
\\[2pt]
\textit{4 bytes}
\end{center}

Não vamos entrar em detalhes sobre representação de números negativos (complemento de dois) neste momento — isso terá seu espaço mais à frente. O que importa agora é fixar uma cadeia de raciocínio que vai acompanhar todo o restante da apostila:

\begin{center}
\textbf{variável $\rightarrow$ tipo $\rightarrow$ bytes $\rightarrow$ memória.}
\end{center}

Toda variável que você declara reserva um espaço de memória cujo tamanho é determinado pelo seu tipo. Essa ideia é a base sobre a qual construiremos, mais adiante, os capítulos de ponteiros e de \textit{structs} — então vale a pena parar um instante e garantir que essa imagem mental ficou clara antes de avançar.

\section{Constantes e Modificadores}
\label{sec:constantes-modificadores}
Às vezes queremos declarar uma variável cujo valor não deve mudar ao longo do programa. Para isso, usamos a palavra-chave \texttt{const}:

\begin{lstlisting}
const int idade = 20;
\end{lstlisting}

Qualquer tentativa de modificar \texttt{idade} depois disso resultará em erro de compilação — o que é útil para proteger valores que representam configurações fixas, limites ou constantes matemáticas.

C também oferece \textbf{modificadores de tipo}, que ajustam o tamanho ou a forma como um número é interpretado:

\begin{lstlisting}
unsigned int  // apenas valores não negativos
long int      // inteiro com mais bits (maior faixa de valores)
short int     // inteiro com menos bits (faixa menor)
\end{lstlisting}

Não vamos nos aprofundar nesses modificadores agora pois eles terão uma seção própria mais adiante, quando discutirmos faixas de valores e \textit{overflow} com mais detalhe.

\section{Operadores}
Operadores são símbolos que indicam ao compilador que operação deve ser realizada sobre um ou mais valores (chamados de \textit{operandos}).

\subsection{Operadores Aritméticos}
\begin{lstlisting}
+   soma
-   subtracao
*   multiplicacao
/   divisao
%   modulo (resto da divisao inteira)
\end{lstlisting}

\subsection{Operadores Relacionais}
Usados para comparar dois valores, sempre resultando em verdadeiro (\texttt{1}) ou falso (\texttt{0}):

\begin{lstlisting}
<    menor que
>    maior que
<=   menor ou igual
>=   maior ou igual
==   igual a
!=   diferente de
\end{lstlisting}

\subsection{Operadores Lógicos}
Usados para combinar expressões relacionais:

\begin{lstlisting}
&&   e (AND)
||   ou (OR)
!    negacao (NOT)
\end{lstlisting}

\subsection{Operadores de Atribuição}
\begin{lstlisting}
=    atribuicao simples
+=   soma e atribui
-=   subtrai e atribui
*=   multiplica e atribui
/=   divide e atribui
%=   calcula o modulo e atribui
\end{lstlisting}

Por exemplo, \texttt{total += 5;} é equivalente a \texttt{total = total + 5;}.

\subsection{Incremento e Decremento}
\begin{lstlisting}
++   incremento
--   decremento
\end{lstlisting}

\begin{lstlisting}
i++;
++i;
\end{lstlisting}

Vale destacar algo que costuma passar despercebido: \textbf{\texttt{++} é um operador, e não simplesmente uma abreviação textual de \texttt{i = i + 1}}. Essa distinção parece sutil agora, mas vai se tornar muito relevante quando estudarmos ponteiros e a diferença entre as formas pré-fixada (\texttt{++i}) e pós-fixada (\texttt{i++}) em expressões mais complexas, como \texttt{vetor[i++]}.

\section{Conversão de Tipos}
Em C, é possível — e às vezes necessário — converter um valor de um tipo para outro. Essa conversão é chamada de \textit{casting}.

\begin{lstlisting}
int a = 5;
int b = 2;

double resultado = (double) a / b;
\end{lstlisting}

Aqui está uma diferença importante em relação ao Python: em Python, \texttt{5 / 2} já retorna \texttt{2.5} automaticamente. Em C, \texttt{a / b} — sendo ambos \texttt{int} — realiza uma \textbf{divisão inteira}, descartando a parte decimal e resultando em \texttt{2}. Para obter o resultado com casas decimais, é preciso converter explicitamente um dos operandos para \texttt{double} (ou \texttt{float}) antes da divisão, como feito no exemplo acima.

Conversões também acontecem de forma implícita, quando atribuímos um valor de um tipo a uma variável de outro tipo:

\begin{lstlisting}
int x = 10;
double y = x;      // conversao implicita: int -> double (sem perda de dados)
\end{lstlisting}

\begin{lstlisting}
double x = 10.5;
int y = (int) x;    // conversao explicita: double -> int (perde a parte decimal, y = 10)
\end{lstlisting}

De forma geral, converter de um tipo "menor" para um "maior" (como \texttt{int} para \texttt{double}) é seguro. O caminho inverso pode causar perda de informação, e por isso é uma boa prática torná-lo explícito com um \textit{cast}, como em \texttt{(int) x}.

\section{Entrada e Saída de Dados}
No Capítulo 1, você já usou a função \texttt{printf} para exibir uma mensagem fixa na tela. Agora que conhecemos variáveis e tipos, podemos usá-la — e sua "parceira" \texttt{scanf} — de forma sistemática, para exibir e ler valores dinâmicos.

\begin{lstlisting}
printf("%d", idade);
scanf("%d", &idade);
\end{lstlisting}

Cada tipo de dado tem um \textbf{especificador de formato} correspondente, usado tanto em \texttt{printf} quanto em \texttt{scanf}:

\begin{lstlisting}
%d   int
%f   float
%lf  double
%c   char
%s   string (vetor de char)
\end{lstlisting}

\begin{errocomum}
\begin{lstlisting}
scanf("%d", idade);   // errado
scanf("%d", &idade);  // correto
\end{lstlisting}
O \texttt{scanf} precisa saber \textbf{onde}, na memória, deve escrever o valor lido — e não apenas qual valor a variável tem no momento. Por isso, é necessário passar o \textbf{endereço de memória} da variável, usando o operador \texttt{\&} antes do nome dela. Omitir o \texttt{\&} é um dos erros mais comuns entre quem está começando em C, e o compilador nem sempre vai te avisar com clareza sobre isso.
\end{errocomum}

Isso naturalmente levanta uma pergunta:

\begin{quote}
\textit{"Por que preciso colocar \texttt{\&}?"}
\end{quote}

A resposta completa para essa pergunta pertence ao capítulo sobre ponteiros, mais adiante. Por ora, fica o registro consciente de que existe algo acontecendo por trás dessa sintaxe que ainda vamos explorar em profundidade — é importante não esconder isso do aluno, mas também não se aprofundar antes da hora.

\section{Escopo e Tempo de Vida}
Toda variável tem um \textbf{escopo} — a região do código onde ela é visível e pode ser usada — e um \textbf{tempo de vida} — o período durante o qual ela existe na memória.

\begin{itemize}
    \item \textbf{Escopo local:} variáveis declaradas dentro de uma função (ou de um bloco delimitado por chaves) só existem e só podem ser acessadas dentro daquele bloco.
    \item \textbf{Escopo global:} variáveis declaradas fora de qualquer função, no topo do arquivo, podem ser acessadas por qualquer função do programa.
\end{itemize}

\begin{lstlisting}
int contador_global = 0; // escopo global

int main(void) {
    // escopo local, existe apenas dentro de main
    int idade_local = 20; 
    return 0;
}
\end{lstlisting}

De forma geral, variáveis locais existem apenas enquanto a função em que foram declaradas está em execução — assim que a função termina, esse espaço de memória é liberado. Variáveis globais, por outro lado, existem durante toda a execução do programa.


\section*{Resumo do Capítulo}
Neste capítulo, deixamos de tratar variáveis como simples "caixas de valores" e passamos a enxergá-las como regiões de memória com tipo, tamanho e representação bem definidos. Vimos os tipos fundamentais de C, como declarar e inicializar variáveis, os principais operadores da linguagem, como converter entre tipos e como realizar entrada e saída básica de dados. Também demos os primeiros passos sobre escopo — um assunto que vai se tornar cada vez mais importante conforme nossos programas crescerem.

\begin{objetivos}
\begin{itemize}
    \item Declarar e inicializar variáveis em C.
    \item Identificar os principais tipos fundamentais.
    \item Explicar a relação entre tipos e representação dos dados na memória.
    \item Utilizar operadores aritméticos, relacionais e lógicos.
    \item Realizar conversões de tipos.
    \item Utilizar \texttt{printf} e \texttt{scanf} para entrada e saída simples.
    \item Diferenciar escopo local e global.
\end{itemize}
\end{objetivos}

\section*{Exercícios}

\begin{exercicio}

\subsection*{Fixação}
Declare variáveis de tipos \texttt{int}, \texttt{float} e \texttt{char} para representar, respectivamente, a idade, a altura (em metros) e a inicial do primeiro nome de uma pessoa. Inicialize cada uma com um valor de sua escolha e utilize \texttt{printf} para exibir os três valores na tela, cada um com o especificador de formato adequado.


\subsection*{Aplicação}
Escreva um programa que leia, via \texttt{scanf}, dois números inteiros digitados pelo usuário. O programa deve calcular e exibir a soma, a diferença, o produto e o resto da divisão entre eles.


\subsection*{Desafio}
Analise o programa a seguir:

\begin{lstlisting}
int a = 5;
int b = 2;

printf("%d\n", a / b);
printf("%f\n", (double)a / b);
\end{lstlisting}

Execute mentalmente (ou compile e rode, se preferir) esse trecho e responda: \textbf{por que as duas operações produzem resultados diferentes?}

Tente formular sua resposta usando os conceitos de tipo e representação de dados discutidos neste capítulo — e não apenas como uma regra decorada sobre operadores.


\end{exercicio}



\begin{tcolorbox}[title=Material Complementar]

Neste capítulo foram utilizados diversos exemplos práticos. As implementações completas, listas adicionais de exercícios, slides da monitoria e perguntas frequentes encontram-se disponíveis no repositório da disciplina.

\end{tcolorbox}



%----------------------------------------
